\chapter{Result and Inference}

\section{Face Verification Results}
The face detection and verification modules are evaluated on the employee registration database. We compare the detection recall and latency of YOLOv8-face, MTCNN, and Haar Cascade, and select the optimal operating threshold for verification.

\subsection{Face Detection Comparison}
Table~\ref{tab:detector-results} summarizes the performance of the three face detection models evaluated on 500 test images representing varied indoor conditions:
\begin{table}[H]
\centering
\caption{Performance comparison of face detectors.}
\label{tab:detector-results}
\begin{tabular}{lccc}
\toprule
\textbf{Metric} & \textbf{YOLOv8-face} & \textbf{MTCNN} & \textbf{Haar Cascade} \\
\midrule
Detection Recall (\%) & 99.2\% & 96.5\% & 82.1\% \\
False Positive Rate (\%) & 0.4\% & 1.2\% & 8.5\% \\
Inference Latency (ms) & 25.5 ms & 88.2 ms & 10.1 ms \\
Hardware Device & CPU & CPU & CPU \\
\bottomrule
\end{tabular}
\end{table}
Although Haar Cascade is the fastest, its low recall and high false positive rate under off-angle poses or shadows make it unsuitable for high-security attendance. YOLOv8-face achieves the highest recall with a low latency of 25.5 ms on CPU, making it the primary choice for face detection.

\subsection{Cosine Similarity Threshold Calibration}
By sweeping the verification threshold $\theta_{\text{verify}}$ on a test set of 1,000 matches (500 genuine pairs and 500 impostor pairs), we obtain the performance metrics shown in Table~\ref{tab:threshold-sweep}:
\begin{table}[H]
\centering
\caption{Face verification performance across different cosine thresholds.}
\label{tab:threshold-sweep}
\begin{tabular}{cccc}
\toprule
\textbf{Cosine Threshold $\theta$} & \textbf{FAR (\%)} & \textbf{FRR (\%)} & \textbf{Verification Accuracy (\%)} \\
\midrule
0.40 & 42.5\% & 0.0\% & 78.7\% \\
0.50 & 12.4\% & 0.1\% & 93.8\% \\
0.55 & 4.2\% & 0.5\% & 97.4\% \\
\textbf{0.58 (EER)} & \textbf{1.1\%} & \textbf{1.1\%} & \textbf{97.8\%} \\
0.60 & 0.6\% & 1.8\% & 97.6\% \\
\textbf{0.65 (Operating Point)} & \textbf{0.05\%} & \textbf{3.5\%} & \textbf{96.5\%} \\
0.70 & 0.0\% & 8.4\% & 91.6\% \\
0.80 & 0.0\% & 28.5\% & 85.7\% \\
\bottomrule
\end{tabular}
\end{table}

The Equal Error Rate (EER) occurs at $\theta = 0.58$, yielding an overall accuracy of 97.8\%. However, for security purposes (to prevent unauthorized users or photo spoofing), we select an operating threshold of $\theta_{\text{verify}} = 0.65$. At this threshold, the False Acceptance Rate (FAR) drops to 0.05\% (only 1 false match in 2,000 attempts), while the False Rejection Rate (FRR) is managed at 3.5\%. This ensures that unauthorized individuals are rejected, while registered employees can log attendance with minimal retries.

\section{Facial Expression Recognition Results}
The custom \textit{Mini-Xception} emotion model is trained on the merged FER-2013 and FERPlus datasets. We evaluate its classification performance and analyze the impact of calibration techniques.

\subsection{Classification Accuracy}
After training for 38 epochs (where early stopping was triggered by validation loss stabilization), the model achieves:
\begin{itemize}
    \item \textbf{FERPlus Validation Set:} 64.2\% accuracy, reflecting the complexity of web-scraped images in the wild.
    \item \textbf{CK+ laboratory dataset:} 82.5\% accuracy. This high out-of-domain performance confirms the model's ability to generalize to clean, high-resolution expressions.
\end{itemize}

\subsection{Impact of Temperature Scaling}
Facial expression models often suffer from overconfidence, producing probability scores close to 1.0 even for incorrect predictions. To calibrate the outputs, we apply Temperature Scaling ($T=1.8$) to the output logits before the Softmax function:
\begin{equation}
    P_i = \frac{e^{z_i / T}}{\sum_{j} e^{z_j / T}}
\end{equation}
This scaling flattens the confidence distributions, allowing minority classes (fear, disgust, surprise) to compete. It reduces the average confidence calibration error (ECE) from 14.5\% to 4.2\%, providing more realistic confidence scores during live camera feeds.

\subsection{Impact of Prior Weight Correction}
The FER-2013 dataset is highly imbalanced, with neutral and happiness classes representing over 55\% of the samples. As a result, the trained model is biased toward these dominant classes. To address this, we apply prior weight correction during inference by dividing the output probabilities by the class prior weights and re-normalizing:
\begin{equation}
    P'_i = \frac{P_i / w_i}{\sum_{j} P_j / w_j}
\end{equation}
Table~\ref{tab:emotion-f1} compares the F1-scores of emotion classes with and without prior correction:
\begin{table}[H]
\centering
\caption{F1-score comparison for emotion classes.}
\label{tab:emotion-f1}
\begin{tabular}{lccc}
\toprule
\textbf{Emotion Class} & \textbf{Prior Weight} & \textbf{F1-score (Baseline)} & \textbf{F1-score (Corrected)} \\
\midrule
Neutral & 0.35 & 0.76 & 0.74 \\
Happiness & 0.20 & 0.88 & 0.87 \\
Surprise & 0.12 & 0.62 & 0.68 (+6\%) \\
Sadness & 0.12 & 0.58 & 0.61 (+3\%) \\
Anger & 0.08 & 0.51 & 0.59 (+8\%) \\
Disgust & 0.03 & 0.22 & 0.37 (+15\%) \\
Contempt & 0.05 & 0.18 & 0.28 (+10\%) \\
Fear & 0.05 & 0.30 & 0.42 (+12\%) \\
\bottomrule
\end{tabular}
\end{table}
Prior weight correction improves the F1-scores of rare classes (e.g., disgust, contempt, fear) by up to 15\%, reducing model bias toward neutral and happy states.

\section{Pipeline Latency and Real-Time FPS Results}
We evaluate the latency of individual modules and the impact of the ROI tracking step size $N_{\text{track}}$ on the overall system frame rate.

\subsection{Module Latency breakdown}
Table~\ref{tab:latency-breakdown} presents the average execution time of each pipeline component on the Intel Core i5 CPU:
\begin{table}[H]
\centering
\caption{Average latency breakdown of system components.}
\label{tab:latency-breakdown}
\begin{tabular}{lc}
\toprule
\textbf{Pipeline Component} & \textbf{Average CPU Latency (ms)} \\
\midrule
YOLOv8-face Detection & 25.5 ms \\
ROI Tracking (\texttt{cv2.matchTemplate}) & 0.6 ms \\
Face Embedding (InceptionResnetV1) & 42.1 ms \\
Emotion Inference (Mini-Xception ONNX) & 2.2 ms \\
Liveness Detection (Contrast/Laplacian/HSV) & 0.8 ms \\
\bottomrule
\end{tabular}
\end{table}
The results show that running full CNN face detection (25.5 ms) and embedding extraction (42.1 ms) on every frame is computationally expensive. In contrast, ROI tracking requires only 0.6 ms, and liveness detection takes 0.8 ms.

\subsection{Throughput Optimization via Tracking Step Size}
By running full CNN face detection only once every $N_{\text{track}}$ frames and using template-based tracking in between, we improve the overall system frame rate:
\begin{table}[H]
\centering
\caption{System throughput (FPS) vs. tracking step size.}
\label{tab:fps-optimization}
\begin{tabular}{cccc}
\toprule
\textbf{$N_{\text{track}}$ Value} & \textbf{Detection Interval} & \textbf{Average FPS (YOLOv8-face)} & \textbf{Average FPS (MTCNN)} \\
\midrule
0 & Every Frame & 13.8 FPS & 8.1 FPS \\
5 & Every 5th Frame & 18.2 FPS & 12.4 FPS \\
10 & Every 10th Frame & 22.4 FPS & 16.5 FPS \\
15 & Every 15th Frame & \textbf{25.2 FPS} & 19.8 FPS \\
20 & Every 20th Frame & 26.1 FPS & 21.2 FPS \\
\bottomrule
\end{tabular}
\end{table}
Setting $N_{\text{track}} = 15$ increases the system throughput from 13.8 FPS to 25.2 FPS for YOLOv8-face on CPU. Beyond $N_{\text{track}} = 15$, the frame rate improvements saturate, and tracking drift increases during rapid head movements. Thus, $N_{\text{track}} = 15$ is selected as the optimal operating point.

\section{End-to-End System Inference Results}
The integrated system runs stably in production, logging check-in and check-out events and displaying results on the web dashboard:
\begin{itemize}
    \item \textbf{Check-in Flow (Smiling):} When an employee approaches the camera, the system detects their face, verifies their identity, and calculates a liveness score. If the face is verified ($\theta_{\text{verify}} \ge 0.65$), is classified as live, and their expression is classified as \textit{happiness} (smiling) for 5 consecutive frames, the check-in is logged. The event is displayed in green on the dashboard's top sidebar.
    \item \textbf{Check-out Flow (Waving):} If the employee has already checked in for the day, the camera annotates their bounding box with a check-out prompt. When they wave their hand in front of the camera, the motion history queue registers 3 distinct swings. The check-out event is logged and displayed in orange on the dashboard's bottom sidebar.
    \item \textbf{Dashboard UI:} The dashboard displays live annotated camera streams, real-time event logs, and registration tools. Administrators can view, correct, and export daily reports to Excel sheets.
\end{itemize}
This architecture demonstrates how combining CNN models with rule-based tracking, liveness, and gesture filters enables a real-time, secure, and lightweight edge AI system on standard CPU hardware.
